\taughtsession{Lecture}{Application Layer Security}{2026-05-12}{09:00}{Fahad}{}

There is information on Moodle relating to the final exam.

\section{Needham Schroeder Symmetric Key Protocol}
The \textit{Needham Schroder Symmetric Key} (NSSK) protocol is now obsolete but is a theoretical underpinning concept for Kerberos, talked about later in this lecture.

NSSK uses a \textit{Key Distribution Centre} (KDC) to distribute shared session keys using only symmetric key cryptography. This works through each party, $A$ and $B$ independently holding a shared symmetric key with the KDC, $S$ in advance. $A$ holds $K_{as}$ and $B$ holds $K_{bs}$. 

Then when $A$ and $B$ want to communicate, the KDC helps establish a fresh shared key between them, $K_{ab}$. This is generated using single-use random number (Nonces), $N_a$ and $N_b$, which are used to help authenticate the parties. The message flow can be seen below:

\begin{enumerate}
    \item $A \rightarrow S$: $A$, $B$, $N_a$\\
    $A$ requests $S$ to supply a key for communication with $B$.
    \item $S \rightarrow A$: $\{N_a, B, K_{ab}, \{K_{ab}, A\}_{K_{bs}}\}_{K_{as}}$\\
    $S$ returns a message encrypted in $A$'s secret key ($K_a$) containing the newly generated $K_{ab}$ and a ticket encrypted with $B$'s secret key. The nonce $N_a$ demonstrates that the message was sent in response to the preceding one. $A$ believes that $S$ sent the message because only $S$ knows $A$'s secret key.
    \item $A \rightarrow B$: $\{K_{ab}, A\}_{K_{bs}}$\\
    $A$ sends the ticket to $B$.
    \item $B \rightarrow A$: $\{N_b\}_{K_{ab}}$\\
    $B$ decrypts the ticket and uses the new key, $K_{ab}$ to encrypt another nonce, $N_b$. 
    \item $A \rightarrow B$: $\{N_b - 1\}_{K_{ab}}$\\
    $A$ demonstrates to $B$ that it was the sender of the previous message by returning an agreed transformation of $N_b$. 
\end{enumerate}

There are, however, weaknesses with NSSK. Message 3 is weak, in that $B$ has no reason to believe the message is fresh; there's no nonce to prove this. If an intruder was to obtain $K_{ab}$ and $\{K_{ab}, A\}_{K_{bs}}$ then the intruder can initiate in a subsequent exchange with $B$ impersonating $A$. For this attack to occur, an old value of $K_{ab}$ has to have been compromised. 

This weakness can be remediated by adding a nonce (timestamp) to message 3. Message 3 would then become $\{K_{ab}, A, t\}_{K_{bs}}$. $B$ decrypts this message and checks that $t$ is recent. This is the solution that has been adopted in Kerberos.

\section{Kerberos}
\textit{Yup, this has been covered in IT \& Internetworking Security, and will be repeated here. Both slightly differently.}

\textit{Kerberos} was developed at MIT in the 1980s to provide authentication and security services for the MIT campus computing network. It has gone through a number of revisions and is widely adopted. In this lecture, we'll consider Kerberos version 5 as specified in RFC 4120. 

Kerberos involves three security objects:
\begin{define}
    \item[Ticket] A token issued to a client by the \textit{Kerberos Key Distribution Centre} (KDC) to present to a particular server, verifying that the sender has recently been authenticated by Kerberos. The ticket includes an expiry time and a newly granted session key for use by the client and the server.
    \item[Authenticator] A token constructed by a client and sent to a server to prove the identity of the user, and also is the currency of any communication with a server. An authenticator can be used only once. It contains the client's name and a timestamp, and is encrypted using the appropriate session key.
    \item[Session Key] A secret key randomly generated by Kerberos and issued to a client for use when communicating with a particular server. Encryption is not mandatory for all communication with servers; the session key is used for encrypting communication with those servers that demand it, and for encrypting all authenticators. 
\end{define}

A client must possess a ticket and session key for each server that they use. The tickets are granted to clients with a lifetime of several hours (at the cyber security team's discretion) so that they can be used for interaction with a particular server until they expire, at which point re-authentication is required. 

Each \textit{Key Distribution Centre} (KDC) contains within it an \textit{Authentication Service} (AS) and a \textit{Ticket Granting Service} (TGS). On login, users are authenticated by the AS. The client is supplied with a TGS ticket and a session key for subsequent communication with the TGS. The client can use the TGS ticket to obtain tickets and session keys for specific services, called \textit{server keys}, from the TGS.

\subsection{Kerberos' Relationship to NSSK}
As alluded to above, Kerberos is related to NSSK. Kerberos follows NSSK quite closely, however Kerberos adds more Nonces which serve two purposes:
\begin{itemize}
    \item To guard against replays of old messages intercepted in the network or the reuse of old tickets found in the memory of machines from which the authorised user has logged out; and
    \item To apply a lifetime to tickets, enabling the system to revoke user's rights if they cease to be authorised users of the system.
\end{itemize}

\subsection{Authentication Process}
\begin{extlink}
There's a diagram of the following Kerberos process in the slides on Moodle.    
\end{extlink}

The below steps show how Kerberos works:

\textbf{Part A: Obtain Kerberos Session key \& TGS ticket, happens once per login session}
\begin{enumerate}
    \item $C \rightarrow A$: $C, T, n$\\
    Client $C$ requests the Kerberos authentication server $A$ to supply a ticket for communicating with the ticket-granting service $T$. 
    \item $A \rightarrow C$: $\{K_{ct}, n\}_{K_{c}}, \{ticket(C,T)\}_{K_{t}}$\\
    $A$ returns a message containing a ticket encrypted in it's secret key and a session key for $C$ to use with $T$. The inclusion of nonce $n$ encrypted in $K_c$ shows that the message comes from the recipient of message 1, who must know $K_c$. The ticket contains $C, T, t_1, t_2, K_{ct}$ where $t_1$ and $t_2$ are the start and end time for the ticket to be valid.
\end{enumerate}

\textbf{Part B: Obtain ticket for a server, once per client-server session}
\begin{enumerate}
    \setcounter{enumi}{2}
    \item $C \rightarrow T$: $\{auth(C)\}_{K_{ct}}, \{ticket(C,T)\}_{K_{t}}, S, n$\\
    $C$ requests the ticket granting server, $T$, to supply a ticket for communication with another server, $S$.
    \item $T \rightarrow C$: $\{K_{cs}, n\}_{K_{ct}}, \{ticket(C,S)\}_{K_{s}}$\\
    $T$ generates the ticket. if it is valid, $T$ generates a new random session key, $K_{cs}$, and returns it with a ticket for $S$ (encrypted in the server's secret key, $K_s$). 
\end{enumerate}

\textbf{Part C: Issue a server request with a ticket}
\begin{enumerate}
    \setcounter{enumi}{4}
    \item $C \rightarrow T$: $\{auth(C)\}_{K_{cs}}, \{ticket(C,S)\}_{K_{s}}, request, n$\\
    $C$ sends the ticket to $S$ with a newly generated authenticator for $C$ and a request. The request would be encrypted in $K_{cs}$ if secrecy of the data is required.
\end{enumerate}

\textbf{Part D: Client authenticates server [optional]}
\begin{enumerate}
    \setcounter{enumi}{5}
    \item $S \rightarrow C$: $\{n\}_{K_{cs}}$\\
    $S$ sends the nonce to $C$, encrypted in $K_{cs}$. 
\end{enumerate}

\subsection{Using Kerberos to Login}
When a user logs into a workstation, the login program sends the user's name to the Kerberos \textit{Authentication Service} (AS). If the user is known to the AS, it replies with a session key, a nonce encrypted in a shared key derived from the user's hashed password, and a ticket for the \textit{Ticket Granting Service} (TGS).

The login program then attempts to decrypt the session key and the nonce using a key derived from the password that the user typed in response to the password prompt. If the password is correct, the login program obtains the session key and the nonce. It checks the nonce and stores the session key with the ticket for subsequent use when communicating with the TGS.

The login program can erase the user's password from its memory, since the ticket now serves to authenticate the user. A login session is the started for the user on the workstation.

Note that the user's password is never exposed to eavesdropping on the network, it is retained wholly within the workstation. 

\subsection{Accessing Servers with Kerberos}
Whenever a program running on a workstation needs to access a new service, it requests a ticket for the service from the TGS. 

When a UNIX user wishes to log into a remote computer, the \textit{rlogin} command prompt on the user's workstation obtains a TGS ticket for access to the \textit{rlogind} network service. The \textit{rlogin} command program sends the TGS ticket, together with a new authenticator, in a request to the \textit{rlogind} process on the computer where the user wishes to login. The \textit{rlogind} program decrypts the TGS ticet with the \textit{rlogin} service's secret key and checks its validity. The \textit{rlogind} program then uses the session key included in the ticket to decrypt the authenticator and checks it is fresh.

There is no need to check the user's name and password following these checks, as the user's identity is known to the \textit{rlogind} program and a login session is established for that user on the remote machine.

\section{Access Control Policies}
Access privileges are specified and subjects' access to objects are determined through a security policy.
\begin{define}
    \item[DAC] Discretionary Access Control Policy - controls access based upon identity
    \item[MAC] Mandatory Access Control Policy - controls access based upon security labels
    \item[RBAC] Role-Based Access Control Policy - controls access based on roles
    \item[ABAC] Attribute Based Access Control Policy - controls access based on attributes   
\end{define}

\section{SAML}
The OASIS \textit{Security Association Markup Language} (SAML) standard defines an XML-based framework for describing and exchanging security information between online business partners. This security information is expressed in the form of portable SAML assertions that applications working across the boundaries of security domains can trust. The OASIS SAML standard defines precise syntax and rules for requesting, creating, communicating and using these SAML assertions. SAML is behind many `single sign on' applications.

\subsection{Single Sign On}
\textit{Single Sign On} (SSO) is the idea of logging in once and this identity being usable across multiple different applications, websites, or service. 

The single identity which exists across multiple applications is known as a \textit{federated identity} and it will be based on a business arrangement. The \textit{identity provider} (IdP) asserts to the \textit{service provider} (SP) that:
\begin{itemize}
    \item the user is known (by referring to the user by their federated identity)
    \item has authenticated to it
    \item and has certain identity attributes
\end{itemize}

SAML 2.0 supports numerous variations on the two main flows (as below) that deal with requirements for:
\begin{itemize}
    \item using various types and strengths of user authentication methods
    \item alternative formats for expressing federated identities
    \item use of different bindings for transporting the protocol messages
    \item inclusion of identity attributes
\end{itemize}

The two main flows are shown below.

\subsubsection{SP Initiated SSO}
SP initiated SSO begins with the user navigating to an SP resource they need to access. As the user is not currently logged in, the SP notices this and sends the user to an IdP to authenticate. The IdP builds and assertion representing the user's authentication at the IdP and sends the user back to the SP with the assertion. The SP processes the assertion and determines whether to grant the user access to the resource.

\subsubsection{IdP Initiated SSO}
In this flow, the user has already authenticated to an IdP and they navigate to a partner SP. The IdP builds an assertion representing the user's authentication state at the IdP and sends the user's browser over to the SP's assertion consumer service. The SP processes the assertion and creates a local security context for the user at the SP. This requires the IdP to be configured with links to partner SP sites.

\subsection{Federated Identity}
SAML can dynamically establish a federated identity during web SSO. Local identities must be linked to the federated identity used to represent the user when interacting with a partner.

The process of associating a federated identifier with the local identity at a partner is called account linking.

\begin{example}{Federated Identity Management \& Account Linking}
In this example there is one IdP: airline.example.com, and two SPs: cars.example.com and hotels.example.com. The user, John Doe, is registered on all three provider sites but the local accounts all have different account identifiers \textit{johndoe}, \textit{jdoe}, and \textit{johnd} respectively.

The sites have agreed to use persistent privacy-preserving pseudonyms for the user's federated identities. John hasn't previously federated identities between these sites.

John starts off his very expensive endeavour by booking flights at airline.example.com. He uses his \textit{johndoe} account.

John then visits cars.example.com to reserve a car. The site sees the user isn't logged in locally, but has previously visited an IdP partner site (airline.example.com). So cars.example.com asks John for consent to federate identities. John consents and his browser redirects back to airline.example.com where the site creates a new pseudonym (\textit{azqu3H7}) for John's use when he visits cars.example.com. The pseudonym is linked to his \textit{johndoe} account. 

John is then redirected to cars.example.com with a SAML assertion indicating user \textit{azqu3H7} is logged in at the IdP. As this is the first time that cars.example.com has seen this identifier, it doesn't know which local account it refers to. John must login at cars.example.com using his \textit{jdoe} account. Then cars.example.com attaches the identity \textit{azqu3H7} to the local \textit{jdoe} account for future use with the IdP. 

After reserving a car, John heads to hotels.example.com in order to book a hotel room. The federation process is repeated with the IdP airline.example.com, creating a new pseudonym (\textit{f78q9C0}) for IdP user \textit{johndoe} that will be used when visiting hotels.example.com. 

John is redirected back to the hotels.example.com SP with a new SAML assertion. The SP requires John to log into his local \textit{johnd} user account and adds the pseudonym as the federated name identifier for future use with the IdP airline.example.com. The user accounts at the IdP and this SP are now linked using the federated name identifier \textit{f78q9C0}.

hotels.example.com does not know about the federation with, or John's account on, cars.example.com. Equivalently, cars.example.com does not know about John's account on, or federation with, hotels.example.com. 

In the future, when John needs to make another booking, he will only need to log in once at airline.example.com.

\end{example}